\documentclass[11pt,a4paper]{article}

% ---------- Codifica e lingua ----------
\usepackage[T1]{fontenc}
\usepackage[utf8]{inputenc}
\usepackage[main=italian,provide=*]{babel}

% ---------- Matematica ----------
\usepackage{amsmath,amssymb}
\usepackage{mathtools}

% ---------- Impaginazione ----------
\usepackage[a4paper,margin=2.7cm]{geometry}
\usepackage{caption}
\usepackage{booktabs}
\usepackage{enumitem}
\usepackage{placeins}
\usepackage[colorlinks=true,linkcolor=black,citecolor=black,urlcolor=blue!55!black]{hyperref}

\newcommand{\ket}[1]{\lvert #1 \rangle}
\newcommand{\bra}[1]{\langle #1 \rvert}

\title{Parte 2 --- rumore quantistico sul dimero\\
\large Panoramica della pipeline (vista a 2000 piedi)}
\author{Samuele \\ Relatore: Prof.\ Alessandro Chiesa}
\date{}

\begin{document}
\maketitle

\begin{abstract}
\noindent Questo documento è una mappa, non una derivazione: serve a vedere in un colpo
d'occhio come è strutturata la Parte 2 (sistema aperto, dimero) prima di entrare nel
dettaglio di ciascun passo. Scope, canali di rumore e decisioni operative sono quelli
confermati dal relatore (risposta del 2 settembre 2026, \texttt{domande\_relatore.md},
sez.\ 12) e seguono la terminologia delle sue slide (\texttt{7Physical\_implementation.pdf}).
\end{abstract}

\tableofcontents

\section{Da dove partiamo}

La Parte 1 (sistema chiuso) è completa sul dimero: Hamiltoniana esatta, VQE con termine
DM, evoluzione temporale $e^{-iHt}$ via Suzuki--Trotter, e correlazioni dinamiche
$C_{ij}^{\alpha\beta}(t)$ tramite un circuito di tipo Hadamard test. Tutto lavora su uno
stato \emph{puro} $\ket{\psi}$ che evolve unitariamente: non c'è interazione con
l'ambiente, non c'è rumore.

La Parte 2 ripercorre esattamente la stessa catena di stadi --- preparazione dello stato
fondamentale, evoluzione temporale, misura dei correlatori --- ma la accompagna con due
sorgenti di rumore realistiche, quelle di un dispositivo quantistico vero. Non si
introduce fisica nuova nell'Hamiltoniana: si aggiunge un modello di come un computer
quantistico reale esegue (in modo imperfetto) gli stessi circuiti già scritti.

\section{Cosa cambia: da $\ket{\psi}$ a $\rho$}

Il rumore accoppia il sistema a un ambiente che non osserviamo. Uno stato puro non basta
più a descrivere il sistema una volta che tracciamo via i gradi di libertà
dell'ambiente: serve l'operatore densità $\rho$, e l'evoluzione non è più unitaria pura
ma un'operazione quantistica più generale. Il prossimo documento (teoria del Passo 1)
richiama questo formalismo dal punto dove serve; qui basta sapere che da questo momento
in poi si lavora con $\rho$ al posto di $\ket{\psi}$, e che in Qiskit questo corrisponde
a usare \texttt{AerSimulator(method="density\_matrix")} invece del simulatore
statevector.

\section{I passi della pipeline}

\begin{table}[htbp]
\centering
\begin{tabular}{@{}cll@{}}
\toprule
\# & Passo & Cosa cambia rispetto alla Parte 1 \\
\midrule
1 & Modello di rumore & nuovo: nessun analogo in Parte 1 \\
2 & VQE con termine DM, rumoroso & stesso ansatz, simulatore a $\rho$ \\
3 & Quantum simulation (Trotter), rumorosa & stesso circuito, compilato per passo \\
4 & Correlazioni dinamiche, rumorose & stesso Hadamard test, readout incluso \\
5 & Scan sui parametri di rumore & nuovo: nessun analogo in Parte 1 \\
\bottomrule
\end{tabular}
\caption{I cinque passi della Parte 2. I passi 2--4 sono mirror dei corrispondenti
stadi di Parte 1; i passi 1 e 5 sono propri della Parte 2.}
\end{table}
\FloatBarrier

\subsection{Passo 1 --- costruzione del modello di rumore}

Il relatore ha confermato due soli canali (slide 17 del corso): errore di gate
(canale depolarizzante, separato a 1 e 2 qubit) ed errore di lettura (readout).
$T_1$/$T_2$ sono esclusi esplicitamente. Il compito di questo passo è tradurre dati di
calibrazione reali --- pubblicati per un chip esistente --- in un oggetto
\texttt{NoiseModel} di Qiskit, passando dal formalismo a operatori di Kraus richiamato
a lezione. Qui ci si ferma a livello di orientamento: il dettaglio teorico è nel
prossimo documento (\S~\ref{sec:prossimo}).

\subsection{Passo 2 --- VQE con termine DM, sotto rumore}

Si ripete la stima del ground state con lo stesso ansatz PMA-2q.3 già validato in Parte
1 (fedeltà 1 nel caso ideale), ma il circuito viene ora eseguito su un simulatore a
operatore densità con il \texttt{NoiseModel} del Passo 1 agganciato. L'osservabile
d'interesse è quanto l'energia stimata e la fedeltà rispetto al ground state esatto si
allontanano dal caso ideale.

\subsection{Passo 3 --- quantum simulation (Trotter), sotto rumore}
\label{subsec:passo3}

Stesso circuito di Suzuki--Trotter di Parte 1. Qui c'è un'insidia tecnica verificata
esplicitamente: se si transpila \emph{l'intero} circuito a un livello di ottimizzazione
alto, Qiskit riconosce che gli $N$ passi identici formano nel complesso un unico
unitario a 2 qubit e li collassa in un circuito a costo \emph{costante}, indipendente da
$N$. Questo distruggerebbe silenziosamente proprio il compromesso che vogliamo
studiare (più passi $\Rightarrow$ meno errore di Trotter ma più rumore accumulato). La
soluzione è transpilare una volta il singolo passo al suo costo minimo (3 CNOT, già
dimostrato in \texttt{circuito\_compatto\_teoria\_e\_limiti.tex}) e poi ripetere il
blocco già compilato $N$ volte, senza ritranspilare il circuito completo.

\subsection{Passo 4 --- correlazioni dinamiche, sotto rumore}

Stesso circuito Hadamard test di Parte 1 (ancilla + 2 qubit di registro), ora con
rumore su tutti i suoi ingredienti: preparazione dello stato, blocco $U(t)$ ripetuto
(Passo 3), e infine la misura sull'ancilla, che include anche l'errore di lettura del
Passo 1.

\subsection{Passo 5 --- scan sui parametri di rumore}

Il relatore chiede esplicitamente di non fermarsi a un solo punto: si ripete il conto
del Passo 4 (o già del Passo 3) al variare di $(\varepsilon_{1q},\varepsilon_{2q},
p_\text{readout})$, per vedere come si sposta il numero di passi ottimale $N^*$ che
minimizza l'errore totale --- risultato di un compromesso fra errore di Trotter
(diminuisce con $N$) e rumore accumulato (cresce con $N$).

\section{Cosa resta fuori}

Per completezza, dichiarato esplicitamente per non lasciarlo implicito: rilassamento e
defasamento ($T_1$, $T_2$) sono esclusi per decisione del relatore, anche se compaiono
nelle sue slide come canali standard; il readout asimmetrico è opzionale, da fare solo
se resta tempo; le reti neurali (NQS) restano fuori scope come in Parte 1.

\section{Il prossimo documento}
\label{sec:prossimo}

Teoria del Passo 1: dal formalismo a operatore densità e agli operatori di Kraus, al
canale depolarizzante e all'errore di lettura, fino alla conversione fra l'errore medio
di gate $\varepsilon$ (quello riportato nelle calibrazioni pubblicate) e il parametro
$\lambda$ che Qiskit vuole in ingresso.

\end{document}
